% Experiment Template
% 1) Copy this file and rename it, e.g. experiments/experiment_04_firewall_rules.tex
% 2) In main.tex, add:
%    \includeexperiment{XX}{Experiment Title}{experiments/your_new_file_name}
% 3) Keep this file content as sections only (no \experiment command here)

\section{Theory}
Combined transposition cipher layers two transposition techniques---column transposition and rail fence---to increase diffusion. In our implementation, column transposition is applied first, then rail fence. Decryption reverses the order: rail fence first, then column transposition, similar to peeling onion layers.


\section{Algorithm}
\subsection*{Encryption Algorithm}
\begin{enumerate}
  \item Input plaintext \texttt{text} and key length \texttt{k}. Generate a random permutation key of size \texttt{k}.
  \item \textbf{Column transposition:} Pad text to fill a grid of width \texttt{k}, then read columns in permutation order to produce intermediate string.
  \item \textbf{Rail fence:} On the intermediate string, collect even-indexed characters first, then odd-indexed characters, to produce the ciphertext.
  \item Combined call: \texttt{ciphertext = r.encrypt(col.encrypt(text))}
\end{enumerate}

\subsection*{Decryption Algorithm}
\begin{enumerate}
  \item Input ciphertext \texttt{text}, same key length \texttt{k}, and same permutation key.
  \item \textbf{Rail fence (reverse):} Place characters alternately into even and odd positions to recover the intermediate string.
  \item \textbf{Column transposition (reverse):} Fill the grid column-by-column in permutation order, then read row-by-row to recover the padded plaintext.
  \item Combined call: \texttt{plaintext\_padded = col.decrypt(r.decrypt(ciphertext))}
  \item Truncate to original length: \texttt{plaintext = plaintext\_padded.substr(0, orig\_len)}
\end{enumerate}
% \section{Flowchart (motherboard)}
% \begin{figure}[H]
%   \centering
%   \begin{autofitflowchart}
%     \node (start) [terminator] {Start};
%     \node (input) [io, below=of start] {Input};
%     \node (process1) [process, below=of input] {Process};
%     \node (decision1) [decision, below=of process1, yshift=-0.2cm] {Condition?};
%     \node (loopstep) [process, right=of decision1, xshift=1.2cm] {Update / Loop};
%     \node (output) [io, below=of decision1, yshift=-0.3cm] {Output};
%     \node (end) [terminator, below=of output] {End};

%     \draw [line] (start) -- (input);
%     \draw [line] (input) -- (process1);
%     \draw [line] (process1) -- (decision1);

%     \draw [line] (decision1.east) -- node[above]{No} (loopstep.west);
%     \draw [line] (loopstep.north) |- ([xshift=1.0cm]process1.east) |- (process1.east);

%     \draw [line] (decision1) -- node[right]{Yes} (output);
%     \draw [line] (output) -- (end);
%   \end{autofitflowchart}
%   \caption{Flowchart title}
% \end{figure}

\section{Code}
\begin{minted}{cpp}
#include<bits/stdc++.h>
using namespace std;
class column {
  vector<int> key;
public:
  column(int k);
  void showKey();
  string encrypt(string text);
  string decrypt(string text);
};
column::column(int k) {
  for (int i = 0; i < k; i++) {
    key.push_back(i);
  }
  for (int i = 0; i < k*10; i++) {
    int j = rand() % k;
    swap(key[i % k], key[j]);
  }
}
void column::showKey() {
  cout << "\nKey is: ";
  for (int i = 0; i < key.size(); i++)
    cout << key[i] << " ";
  cout << endl;
}
string column::encrypt(string text) {
  int k = key.size();
  int n = text.size();
  int rows = (n + k - 1) / k;
  // pad text with spaces to fill the grid
  text.append(rows * k - n, ' ');
  string encrypted;
  encrypted.reserve(rows * k);
  // read columns in order of key
  for (int colIndex = 0; colIndex < k; ++colIndex) {
    int col = key[colIndex];
    for (int r = 0; r < rows; ++r) {
      encrypted.push_back(text[r * k + col]);
    }
  }
  return encrypted;
}
string column::decrypt(string text) {
  int k = key.size();
  int n = text.size();
  int rows = (n + k - 1) / k;
  // build empty grid
  vector<char> grid(rows * k, ' ');
  int pos = 0;
  // fill columns in order of key
  for (int colIndex = 0; colIndex < k && pos < n; ++colIndex) {
    int col = key[colIndex];
    for (int r = 0; r < rows && pos < n; ++r) {
      grid[r * k + col] = text[pos++];
    }
  }
  // read row-major
  string decrypted;
  decrypted.reserve(rows * k);
  for (int r = 0; r < rows; ++r) {
    for (int c = 0; c < k; ++c) {
      decrypted.push_back(grid[r * k + c]);
    }
  }
  return decrypted;
}
class rail {
  public:
    rail();
    string encrypt(string s);
    string decrypt(string s);
};
rail::rail() {}
string rail::encrypt(string s) {
  string ans;
  int n = s.length();
  if (n%2 == 1) {
    s+=' ';
    n++;
  }
  for (int i = 0; i < n; i+=2)
    ans+= s[i];
  for (int i = 1; i < n; i+=2)
    ans += s[i];
  return ans;
}
string rail::decrypt(string s) {
  string ans = s;
  int n = s.length();
  for (int i = 0; i < n; i++) {
    if (i%2 == 0) {
      ans[i] = s[i/2];
    } else {
      ans[i] = s[n/2 + i/2];
    }
  }
  return ans;
}
class combined {
  private:
    column col;
    rail r;
public:
  combined(int k) : col(k) {}
  void showKey() {
    col.showKey();
  }
  string encrypt(string text) {
    return r.encrypt(col.encrypt(text));
  }
  string decrypt(string text) {
    return col.decrypt(r.decrypt(text));
  }
};
int main() {
  srand(time(0));
  int k;
  cout << "Enter the key length: ";
  cin >> k;
  cin.ignore();
  combined cipher(k);
  cipher.showKey();
  cout << "Enter Plaintext : ";
  string text;
  getline(cin, text);
  int orig_len = text.size();
  string encrypted = cipher.encrypt(text);
  cout << "Encrypted text: " << encrypted << endl;
  string decrypted = cipher.decrypt(encrypted);
  decrypted = decrypted.substr(0, orig_len);
  cout << "Decrypted text: " << decrypted << endl;
  return 0;
}
\end{minted}

\section{Result and Discussion - I/O}
Combined Transposition cipher is quite flexible to use in most cases with very little restriction.
\begin{minted}{text}
Enter the key length: 5

Key is: 1 2 0 4 3
Enter Plaintext : This is encrypted
Encrypted text: hriyTc tspsd  ien e
Decrypted text: This is encrypted
\end{minted}

\section{Conclusion}
We successfully implemented a Combined Transposition Cipher by layering Columnar Transposition and Rail Fence techniques. The experiment demonstrates that composing multiple simple transpositions significantly increases diffusion and obscures plaintext patterns.
